The deployment of deep multi-agent reinforcement learning (MARL) in high-stakes, dynamic environments—such as Real-Time Bidding (RTB) markets and decentralized energy grids—has revealed a fundamental fragility: the inability to maintain robust coordination under non-stationarity \cite{papoudakis2021benchmarking, gronauer2022multi}. In these domains, the environment is not merely a passive stage but a co-evolving system where the shifting strategies of competitors continuously violate the Markov assumptions underlying traditional RL \cite{yuan2023multi}. While communication has long been proposed as a remedy for partial observability \cite{foerster2016learning, sukhbaatar2016learning}, current MARL communication protocols are largely instrumental, optimized solely to maximize a scalar reward. This often leads to ``cheap talk'' or, more perversely, emergent implicit collusion, where agents prioritize joint reward exploitation at the expense of market stability and belief accuracy \cite{hu2021multi, calvo2023survey}.

The gap in current MAS literature is the lack of a **principled communicative objective** that transcends simple utility maximization. Most existing methods treat communication as just another action in the agent's policy space, susceptible to the same variance and instability as physical actions in non-stationary regimes. We argue that for communication to be truly robust, it must be decoupled from immediate instrumental convergence and instead be grounded in the maintenance of a coherent, shared internal model of the world.

In this paper, we introduce a novel framework for multi-agent coordination based on **Active Inference (ActInf)** \cite{friston2010free, parrill2022active}. Active Inference posits that agents act to minimize their Variational Free Energy (VFE), a proxy for the ``surprise'' encountered when their generative models fail to predict sensory input. We extend this principle to communication, framing it as an **epistemic action** \cite{millidge2024sophisticated}. Under our framework, agents communicate not just to achieve a goal, but to actively reduce uncertainty regarding the latent states of the environment and the hidden intentions of other agents. This shift from reward-centric signaling to belief-centric exchange provides a natural regularization against collusion; an agent motivated by belief accuracy is inherently resistant to the deceptive or sub-optimal strategies that often emerge from pure reward-chasing in MARL \cite{smith2022active}.

As illustrated in Figure 1, our approach replaces the traditional reward-driven communication loop with a hierarchical generative model where communication acts as a bridge between the internal models of disparate agents. By minimizing expected free energy, agents select messages that maximize information gain (epistemic value) while simultaneously pursuing long-term goals (extrinsic value). This dual objective allows for the emergence of sophisticated coordination that is significantly more resilient to the ``moving target'' problem of non-stationary environments.

Our core hypothesis is that framing communication as an epistemic requirement for VFE minimization leads to more stable, efficient, and non-collusive coordination than traditional MARL. We validate this through the following contributions:
\begin{itemize}
    \item \textbf{The ActInf-Comm Framework}: A novel hierarchical architecture for multi-agent Active Inference that formalizes communication as an epistemic action within the VFE objective.
    \item \textbf{Collusion Mitigation Theory}: A theoretical demonstration of how the epistemic drive in ActInf inherently penalizes the low-entropy, exploitative equilibria characteristic of emergent collusion in MARL.
    \item \textbf{Empirical SOTA in RTB}: We demonstrate that in simulated high-frequency RTB environments, our framework achieves a $>15\%$ improvement in economic efficiency and robustness over state-of-the-art baselines like QMIX \cite{rashid2018qmix} and MADDPG \cite{lowe2017multi}.
    \item \textbf{Emergent Concept Communication}: We provide qualitative evidence that our agents, augmented by Large Language Models (LLMs) for inference acceleration \cite{wei2022chain}, can communicate novel semantic concepts to adapt to environmental shifts not explicitly encoded in their initial reward functions.
\end{itemize}

The following sections detail the mathematical foundations of our ActInf approach, the implementation of communication as uncertainty reduction, and our comprehensive experimental evaluation in complex, non-stationary market simulations.